% GENERATED FILE. Edit canonical lessons, then run npm run generate:books.
% canonical-source-hash: fnv1a64:14a76f29d6e01c9b
% canonical-lessons: FR-C327-recette, FR-C327-ingredient, FR-C327-pincee, FR-C327-cuilleree, FR-C327-prechauffer

\chapter{La recette}
\label{ch:fr-la-recette}
\hlchaptermodality{327}

\begin{chapteropening}
\textbf{By the end of this chapter:} \emph{I can talk about a recipe: the recipe, the ingredients, a pinch, a spoonful and preheating the oven.}

Complete the last lesson of chapter 327: I can talk about a recipe: the recipe, the ingredients, a pinch, a spoonful and preheating the oven.
\end{chapteropening}

\section[la recette]{la recette — the recipe; the takings}

\label{lesson:FR-C327-recette}



\begin{quote}
Before the new one: say the French for the resistance, then the French for to commemorate.
\end{quote}

\subsection*{You'll want to know: la recette}
\textbf{la recette} — \textquotedblleft{}the recipe; the takings\textquotedblright{}.

\begin{quote}
Ma grand-mère m'a donné sa recette de tarte.
\end{quote}

\begin{quote}
\emph{My grandmother gave me her tart recipe.}
\end{quote}

In a shop, \emph{la recette} is the day's takings.

\begin{cousinweb}
\emph{recette} comes from Old French \emph{recete}, from Latin \emph{recepta} \textquotedblleft{}things received\textquotedblright{} (EtymDB), from \emph{recipere} \textquotedblleft{}to take back, to receive\textquotedblright{}.
\end{cousinweb}

\begin{grammarlens}[title={What you've built}]
One more word for telling, explaining and describing at B1.
\end{grammarlens}

\subsection*{Your turn}
\begin{itemize}
\raggedright
  \item \emph{Say it:} \emph{la recette}
  \item \emph{Say it:} \emph{la recette}, once more
  \item \emph{Say it:} a sentence of your own with \emph{la recette}
  \item \emph{Recall:} say the French for the resistance, then the French for to commemorate, then say \emph{la recette} again
\end{itemize}

\begin{tcolorbox}[breakable,colback=teal!4,colframe=teal!35!black,title={Before you move on}]
What does la recette mean? (\textquotedblleft{}The recipe; the takings\textquotedblright{}.) Say it once more.
\end{tcolorbox}
\section[l'ingrédient]{l'ingrédient — the ingredient}

\label{lesson:FR-C327-ingredient}



\begin{quote}
Before the new one: say the French for to commemorate, then the French for the recipe.
\end{quote}

\subsection*{You'll want to know: l'ingrédient}
\textbf{l'ingrédient} — \textquotedblleft{}the ingredient\textquotedblright{}.

\begin{quote}
Elle a acheté tous les ingrédients au marché.
\end{quote}

\begin{quote}
\emph{She bought all the ingredients at the market.}
\end{quote}

Masculine: \emph{un ingrédient}.

\begin{cousinweb}
\emph{ingrédient} is a learned borrowing from Latin \emph{ingrediēns} \textquotedblleft{}going in\textquotedblright{} (EtymDB), from \emph{ingredī} \textquotedblleft{}to enter\textquotedblright{}, built on \emph{gradī} \textquotedblleft{}to step\textquotedblright{}.
\end{cousinweb}

\begin{grammarlens}[title={What you've built}]
One more word for telling, explaining and describing at B1.
\end{grammarlens}

\subsection*{Your turn}
\begin{itemize}
\raggedright
  \item \emph{Say it:} \emph{l'ingrédient}
  \item \emph{Say it:} \emph{l'ingrédient}, once more
  \item \emph{Say it:} a sentence of your own with \emph{l'ingrédient}
  \item \emph{Recall:} say the French for to commemorate, then the French for the recipe, then say \emph{l'ingrédient} again
\end{itemize}

\begin{tcolorbox}[breakable,colback=teal!4,colframe=teal!35!black,title={Before you move on}]
What does l'ingrédient mean? (\textquotedblleft{}The ingredient\textquotedblright{}.) Say it once more.
\end{tcolorbox}
\section[la pincée]{la pincée — the pinch (of salt, of pepper)}

\label{lesson:FR-C327-pincee}



\begin{quote}
Before the new one: say the French for the recipe, then the French for the ingredient.
\end{quote}

\subsection*{You'll want to know: la pincée}
\textbf{la pincée} — \textquotedblleft{}the pinch (of salt, of pepper)\textquotedblright{}.

\begin{quote}
Il ajoute une pincée de sel.
\end{quote}

\begin{quote}
\emph{He adds a pinch of salt.}
\end{quote}

As much as you hold between two fingers.

\begin{cousinweb}
\emph{pincée} is built on \emph{pincer} \textquotedblleft{}to pinch\textquotedblright{}, whose origin is not settled (EtymDB lists no parent).
\end{cousinweb}

\begin{grammarlens}[title={What you've built}]
One more word for telling, explaining and describing at B1.
\end{grammarlens}

\subsection*{Your turn}
\begin{itemize}
\raggedright
  \item \emph{Say it:} \emph{la pincée}
  \item \emph{Say it:} \emph{la pincée}, once more
  \item \emph{Say it:} a sentence of your own with \emph{la pincée}
  \item \emph{Recall:} say the French for the recipe, then the French for the ingredient, then say \emph{la pincée} again
\end{itemize}

\begin{tcolorbox}[breakable,colback=teal!4,colframe=teal!35!black,title={Before you move on}]
What does la pincée mean? (\textquotedblleft{}The pinch (of salt, of pepper)\textquotedblright{}.) Say it once more.
\end{tcolorbox}
\section[la cuillerée]{la cuillerée — the spoonful}

\label{lesson:FR-C327-cuilleree}



\begin{quote}
Before the new one: say the French for the ingredient, then the French for the pinch.
\end{quote}

\subsection*{You'll want to know: la cuillerée}
\textbf{la cuillerée} — \textquotedblleft{}the spoonful\textquotedblright{}.

\begin{quote}
Elle met deux cuillerées de sucre dans son café.
\end{quote}

\begin{quote}
\emph{She puts two spoonfuls of sugar in her coffee.}
\end{quote}

Also spelt \emph{cuillérée}. \emph{Une cuillerée à soupe} — a tablespoonful.

\begin{cousinweb}
\emph{cuillerée} is built on \emph{cuiller} \textquotedblleft{}spoon\textquotedblright{} (EtymDB), which is inherited from Latin \emph{cochlear} \textquotedblleft{}a spoon\textquotedblright{} (Lewis \& Short).
\end{cousinweb}

\begin{grammarlens}[title={What you've built}]
One more word for telling, explaining and describing at B1.
\end{grammarlens}

\subsection*{Your turn}
\begin{itemize}
\raggedright
  \item \emph{Say it:} \emph{la cuillerée}
  \item \emph{Say it:} \emph{la cuillerée}, once more
  \item \emph{Say it:} a sentence of your own with \emph{la cuillerée}
  \item \emph{Recall:} say the French for the ingredient, then the French for the pinch, then say \emph{la cuillerée} again
\end{itemize}

\begin{tcolorbox}[breakable,colback=teal!4,colframe=teal!35!black,title={Before you move on}]
What does la cuillerée mean? (\textquotedblleft{}The spoonful\textquotedblright{}.) Say it once more.
\end{tcolorbox}
\section[préchauffer]{préchauffer — to preheat}

\label{lesson:FR-C327-prechauffer}



\begin{quote}
Before the new one: say the French for the pinch, then the French for the spoonful.
\end{quote}

\subsection*{You'll want to know: préchauffer}
\textbf{préchauffer} — \textquotedblleft{}to preheat\textquotedblright{}.

\begin{quote}
Il préchauffe le four à 180 degrés.
\end{quote}

\begin{quote}
\emph{He preheats the oven to 180 degrees.}
\end{quote}

\emph{Préchauffer le four} is the first line of many recipes.

\begin{cousinweb}
\emph{préchauffer} is built on \emph{pré-} \textquotedblleft{}before\textquotedblright{} and \emph{chauffer} \textquotedblleft{}to heat\textquotedblright{}, which is inherited from Latin \emph{calefacere} \textquotedblleft{}to make warm\textquotedblright{} (Lewis \& Short).
\end{cousinweb}

\begin{grammarlens}[title={What you've built}]
One more word for telling, explaining and describing at B1.
\end{grammarlens}

\subsection*{Your turn}
\begin{itemize}
\raggedright
  \item \emph{Say it:} \emph{préchauffer}
  \item \emph{Say it:} \emph{préchauffer}, once more
  \item \emph{Say it:} a sentence of your own with \emph{préchauffer}
  \item \emph{Recall:} say the French for the pinch, then the French for the spoonful, then say \emph{préchauffer} again
\end{itemize}

\begin{tcolorbox}[breakable,colback=teal!4,colframe=teal!35!black,title={Before you move on}]
What does préchauffer mean? (\textquotedblleft{}To preheat\textquotedblright{}.) Say it once more.
\end{tcolorbox}
